\phantomsection
\section*{To the Reader}
\addcontentsline{toc}{section}{To the Reader}

Advent opens the Church's year. It begins with the first Evening Prayer of the
Sunday that falls on or nearest 30 November, sung on the Saturday evening, and
it ends on the evening of 24 December, when the first Evening Prayer of
Christmas is sung. In 2026, for example, the First Sunday is 29 November and
Advent ends on the evening of Thursday 24 December. For four weeks, or a little
less, the Church keeps watch for her Lord, and she asks her children to keep it
with her.

She keeps it from three books. The Roman Missal holds the prayers of the Mass;
the dioceses of the United States pray from its third edition, in the English
first used on the First Sunday of Advent in 2011. The Lectionary for Mass holds
the readings. On Sundays it reads in a cycle of three years, called A, B and C, and
Appendix~\ref{sec:daily} gives the letter for the years ahead. On weekdays
Advent reads the same lessons every year. The Liturgy of the Hours
holds the Church's daily prayer: the Office of Readings, with a passage of
Scripture and a second reading from a Father, a saint or a teacher of the
Church; Morning Prayer, which sings the Canticle of Zachary; Evening Prayer,
which sings the Canticle of Mary; and Night Prayer before sleep. All three books
are the Church's own, and a parishioner may take up any of them.

Keep a Bible at hand. The Scripture quoted here is the Douay-Rheims version.
Its numbering sometimes differs from your Lectionary's, most often in the
Psalms; where it does, a reading is named by the Lectionary's numbers and the
quoted words by the Douay's.

A few words recur. At Mass the Entrance Antiphon, called the Introit in the
older books, is the verse sung or said as the priest goes to the altar; the
Collect is the prayer he says in the name of all before the readings; the verse
before the Gospel is sung with the Alleluia, which Advent keeps; and the
Preface is the hymn of thanks with which the Eucharistic Prayer begins. The days
of the calendar have ranks. A solemnity is the highest, and some solemnities
are also holy days of obligation; a feast comes next; a memorial honours a
saint within the day's own celebration, and is either obligatory or optional;
and a saint whose day falls on one of the privileged days before Christmas is
only commemorated, remembered within the Mass and Office of Advent.

The readings of every day are set out in Appendix~\ref{sec:daily}, and the
prayers of the season, in Latin and English, in Appendix~\ref{sec:prayers}. If
you begin late, turn to the day's entry in Appendix~\ref{sec:daily} and start
there.

If this is your first Advent kept with any care, begin small. Go to Mass on the
First Sunday and listen closely to the Gospel; it sets the key for the four
weeks. Learn the Angelus and say it once a day, at morning, noon or evening; its
closing prayer is the prayer of the Fourth Sunday of Advent, and by Christmas
you will know it by heart. Then read one reading a day, the first reading of the
weekday Mass, which until 16 December is almost always from Isaiah. That is
enough to keep Advent. The rest is offered for those who want more, and no one
is asked to do all of it.

When the weeks are over and the night of Christmas comes, the Church will sing
the angels' hymn she has left unsung on the Sundays of Advent. The waiting is
for that night, and for the coming beyond it.
